\sectionguard
\section{Synopses of the Four Dogmas}

Each synopsis condenses the corresponding chapter without erasing its
historical difficulty; the chapter and its notes control. The fields are
uniform so that the four can be compared field by field, and the last two
fields---exclusions and the ecumenical position---are the ones a reader
in a hurry is most likely to need.

\begin{dogmasynopsis}{Divine Motherhood (Theotokos)}
\synfield{Object}{Mary is truly Mother of God: the Son born of her
according to the flesh is the eternal Word in person.}
\synfield{Mode and authority}{Conciliar dogma within Christology:
Ephesus (431) approving Cyril's second letter to Nestorius (DS 250--251);
Chalcedon (451) Definition of Faith; Constantinople II (553) anathemas.}
\synfield{Operative wording}{``\ldots{} they ventured to call the Holy
Virgin, the Mother of God \ldots{} because of her was born that holy body
with a rational soul, to which the Word being personally united is said
to be born according to the flesh'' (Cyril, tr. Percival).}
\synfield{Biblical center}{Lk 1:43 (``mother of my Lord''); Gal 4:4;
Jn 1:14. Scripture supplies the reality, not the word.}
\synfield{Principal early witnesses}{Ignatius (\work{Eph.} 7, 18--19:
``both of Mary and of God''); Irenaeus's Eve--Mary recapitulation
(\work{AH} III.22.4); \work{Sub tuum praesidium} (dating disputed);
Cyril's letters; Formula of Reunion (433); Vincent of L\'erins (c.~434).}
\synfield{Thomistic support or limit}{Fully affirmed; classic locus ST
III, q.~35 (motherhood terminates in the person), cited at that
generality.}
\synfield{Development hinge}{Nestorius's \term{Christotokos} (428)
forced the judgment that motherhood is of the person; the Formula of
Reunion (433) settled the wording both parties could confess.}
\synfield{Doctrinal fruit}{Guards the unity of Christ; ground of all
Marian doctrine; source of the rule that every Marian dogma is first a
Christological one.}
\synfield{Required assent}{Divine and Catholic faith.}
\synfield{Exclusions and open questions}{No origin of divinity in Mary;
no divinization of her; the status of Cyril's twelve anathemas and the
procedural history of Ephesus remain complex and are not part of the
defined object.}
\synfield{Ecumenical position}{Common ground with the Orthodox and with
the classical Reformation confessions; not in dispute among the
communions surveyed in Chapter 9.}
\end{dogmasynopsis}

\begin{dogmasynopsis}{Perpetual Virginity (Aeiparthenos)}
\synfield{Object}{Mary virgin in conceiving, in giving birth, and ever
after: \term{ante partum}, \term{in partu}, \term{post partum}.}
\synfield{Mode and authority}{Ancient dogma of universal Tradition and
the ordinary universal magisterium; creeds; Constantinople II
(``ever-virgin,'' DS 427); exact formulation at Lateran 649, canon 3
(DS 503), papally confirmed synod.}
\synfield{Operative wording}{``\ldots{} \term{semperque Virginem et
immaculatam Mariam \ldots{} incorruptibiliter eum genuisse,
indissolubili permanente et post partum eiusdem virginitate}'' (DS 503).}
\synfield{Biblical center}{Mt 1:18--25 and Lk 1:26--35 (virginal
conception directly); the threefold confession received through
Tradition.}
\synfield{Principal early witnesses}{Creeds and Ignatius (virginal
conception); \work{Protevangelium of James} as reception evidence;
Jerome against Helvidius (383); Leo's Tome (DS 291) on the incorrupt
birth; \term{aeiparthenos} standard from the fifth century.}
\synfield{Thomistic support or limit}{Fully affirmed: ST III, q.~28,
a.~3 gives four reasons against Helvidius, all arguments of fittingness
built on the received dogma.}
\synfield{Development hinge}{Fourth-century controversy (Helvidius,
Jovinian) and the conciliar fixing of \term{aeiparthenos}; canon 3 of 649
is collateral precision inside an anti-Monothelite confession.}
\synfield{Doctrinal fruit}{Sign of God's absolute initiative in the
Incarnation and of Mary's undivided self-gift; image of the Church's
undivided fidelity (LG 63--64).}
\synfield{Required assent}{Divine and Catholic faith.}
\synfield{Exclusions and open questions}{No anatomical theory defined;
no contempt of marriage---and Jerome's own invective against marriage is
not the Church's teaching; the ``brothers'' are identified as kinsmen,
but the Hieronymian (cousins) and Epiphanian (prior marriage) readings
both remain undefined.}
\synfield{Ecumenical position}{Held by the Orthodox; confessed in the
\work{Smalcald Articles} (1537) among articles then not in dispute;
variously held in later Reformation practice.}
\end{dogmasynopsis}

\begin{dogmasynopsis}{Immaculate Conception}
\synfield{Object}{From the first instant of her conception Mary was
preserved free from all stain of original sin, by singular grace, in
view of Christ's merits.}
\synfield{Mode and authority}{Solemn \term{ex cathedra} definition:
Pius IX, \work{Ineffabilis Deus}, 8 December 1854 (DS 2803--2804).}
\synfield{Operative wording}{``\term{in primo instanti suae conceptionis
\ldots{} intuitu meritorum Christi Jesu \ldots{} ab omni originalis
culpae labe praeservatam immunem}'' (DS 2803).}
\synfield{Biblical center}{Gen 3:15 (the enmity, not the disputed
pronoun) and Lk 1:28 (\term{kecharitomene}), read cumulatively within
Tradition, not as isolated proofs.}
\synfield{Principal early witnesses}{New Eve parallel (Justin,
Irenaeus); \term{Panagia} praise; \term{immaculatam} at Lateran 649;
the feast of the Conception, kept in some churches and tolerated at Rome
in Thomas's day.}
\synfield{Thomistic support or limit}{Real limit: Thomas denied
sanctification prior to animation to protect universal redemption (ST
III, q.~27, a.~2), holding that she ``did indeed contract original sin,
but was cleansed therefrom before her birth''; the definition answers his
concern through preservative redemption rather than overruling it.}
\synfield{Development hinge}{Bernard's objection (c.~1140), closing in
submission to Rome; Scotus's preservative-redemption argument
(\work{Ordinatio} III, d.~3, q.~1); papal steering from Sixtus IV
(1476--1483) through Trent (1546) to Alexander VII (1661), whose wording
1854 nearly reproduces.}
\synfield{Doctrinal fruit}{Mary as the most perfect case of redemption;
grace precedes all merit; the Church's own destiny of holiness realized
in one member.}
\synfield{Required assent}{Divine and Catholic faith; formal sanction on
willful denial.}
\synfield{Exclusions and open questions}{Not the virginal conception of
Jesus; personal sinlessness is a distinct teaching; no pre-Nicene witness
states the formula, and the medieval difficulty is honest history; the
1858 Lourdes sentence is private revelation and adds nothing to the
dogma.}
\synfield{Ecumenical position}{Rejected as an innovation by the
Ecumenical Patriarchate's encyclical of 1895; ARCIC (2005) finds the
teaching ``consonant with'' Scripture while leaving required assent
unresolved.}
\end{dogmasynopsis}

\begin{dogmasynopsis}{Assumption}
\synfield{Object}{Mary, her earthly life completed, was assumed body
and soul into heavenly glory.}
\synfield{Mode and authority}{Solemn \term{ex cathedra} definition:
Pius XII, \work{Munificentissimus Deus} \S44, 1 November 1950 (DS
3903--3904), after the 1946 consultation of the episcopate.}
\synfield{Operative wording}{``\term{Immaculatam Deiparam semper
Virginem Mariam, expleto terrestris vitae cursu, fuisse corpore et anima
ad caelestem gloriam assumptam}'' (DS 3903).}
\synfield{Biblical center}{No narrative text; canonical horizon: the
Mother beside her Son's lot (\S38), new-Eve victory over sin and death
(Gen 3:15 with 1 Cor 15, \S39); Rev 12 and Ps 132 as reported figures.}
\synfield{Principal early witnesses}{Dormition/Assumption liturgy East
and West (Adrian I's sacramentary, the Gallican books, the Byzantine
office); eighth-century homilists, chiefly John Damascene; the
\work{Transitus} apocrypha as reception only---and listed among rejected
books in the \work{Decretum Gelasianum}; early centuries genuinely
silent.}
\synfield{Thomistic support or limit}{Occasional affirmations only;
``never dealt directly with this question'' (\work{MD} \S31); no
dedicated proof claimed---though Thomas himself argued from the
Assumption when treating Mary's sanctification.}
\synfield{Development hinge}{Scholastic fittingness arguments; the
petitions gathered from 1849 to 1940 and the almost two hundred Fathers
of Vatican I; the 1946 consultation's near-unanimity, read by \work{MD}
\S12 as itself an infallible manifestation.}
\synfield{Doctrinal fruit}{Singular participation in the Resurrection;
anticipation of the Church's own end (CCC 966, 972--974); confirmation of
the hope of bodily resurrection (\work{MD} \S42).}
\synfield{Required assent}{Divine and Catholic faith; formal sanction on
willful denial.}
\synfield{Exclusions and open questions}{Death versus dormition
deliberately left open; no place, date, tomb, or \work{Transitus}
narrative defined; assumption is God's act, not an ascension.}
\synfield{Ecumenical position}{Content largely shared with the Orthodox
Dormition tradition, the mode of definition contested; ARCIC (2005)
affirms the teaching as consonant with Scripture, noting the dogma takes
no position on how Mary's life ended.}
\end{dogmasynopsis}
